\section{Actual arbitrary original section powers in the same ample model}
\label{sec:native-ample-model-powered-sections}

\begin{theorem}[True original section powers give same-model coefficients and norms]
\label{thm:native-ample-model-powered-section-norm}
\lean{EnvelopeProperty.NativeAmpleModelPoweredSections.poweredSection,EnvelopeProperty.NativeAmpleModelPoweredSections.poweredSection_norm}
\leanok
\uses{thm:native-ample-model-section-block-map,thm:scheme-global-section-block-power-norm,thm:scheme-global-section-model-norm-isometry,lem:gl-sections}
Fix genuine original positive-power embedding data identifying the first
pulled projective twist with $L^{\otimes r}$ on the same original integral
scheme. For every entire original section $s$ of $L^{\otimes d}$, take its
actual original $r$th tensor power and transport it through the genuine whole
original exponent isomorphism to $L^{\otimes rd}$. The whole block comparison
sends this actual original section to a section of degree $d$ of the same one
model over the entire original native function field. At every actual original
native point, its genuine model section norm is exactly $\lVert s\rVert^r$.
\end{theorem}
\begin{proof}
The actual original section-power construction lands first in
$L^{\otimes dr}$. The true original exponent equality identifies this whole
module with $L^{\otimes rd}$, and the actual whole section isomorphism preserves
its true model norm. The entire block comparison identifies that same norm
with genuine model section norm under the same whole analytic homeomorphism.
The proved actual original section-power norm formula gives the asserted
$r$th power. No section-power coefficient or metric-comparison premise enters.
\end{proof}

\begin{theorem}[Pointwise nonvanishing and logarithmic values of the same powers]
\label{thm:native-ample-model-powered-section-lognorm}
\lean{EnvelopeProperty.NativeAmpleModelPoweredSections.poweredSection_norm_ne_zero_iff,EnvelopeProperty.NativeAmpleModelPoweredSections.poweredSection_logNorm}
\leanok
\uses{thm:native-ample-model-powered-section-norm}
If $r>0$, the actual powered model section has nonzero norm at a point exactly
when the original whole section has nonzero true norm there. For all $r\ge0$,
its model logarithmic section norm is $r\log\lVert s\rVert$, with the actual
extended-real conventions, including zero original norms.
\end{theorem}
\begin{proof}
Apply the preceding genuine norm identity and the exact nonvanishing criterion
for a positive integer power. For the logarithmic identity, coerce the actual
nonnegative norm into the extended nonnegative reals and apply the genuine
extended-real logarithm-of-power theorem. Zero norms are included by that same
theorem; no positivity assumption on the individual section norm is used.
\end{proof}
